\chapter{XNB Type Readers and Object Graphs}
\label{ch:xnb-type-readers}
% Compatibility label retained from the combined content chapter.
\label{sec:xnb-and-cnj-model}

An XNB file is an object graph, not a bag of independent records. The container selects a table of reader factories; one ContentReader session constructs root and nested objects, records shared-resource fixups, and enforces the bounds described in Chapter~\ref{ch:xnb-container}.

\section{Two registries that must not be confused}

\subsection{Custom \texttt{.xnb} readers use a process-wide registry; loose and \texttt{.cnj} readers do not}

An earlier shorthand in this chapter needs a correction: \texttt{RegisterTypeReader<T>} is
\emph{not} the extension point for a custom \texttt{.xnb}. \texttt{Load<T>()} checks for an
\texttt{.xnb} before consulting its per-manager loose-reader table, and
\texttt{LoadXnbAsset()} gives the binary file's own reader table to the static
\cnaclass{ContentTypeReaderManager}. A custom binary reader instead registers a canonical
reader-name factory in that global registry. CNA has no reflection fallback, whereas FNA first
tries its factory map and can otherwise reflect the named managed type, so this registration is
public and \texttt{CNAEXT}-tagged in CNA rather than FNA's internal AOT fallback.

The registry starts empty on purpose. \texttt{ContentManager}'s constructor does \emph{not}
populate it, and neither does \cnaclass{Game}; otherwise a test that clears the global state could
not exercise the genuine ``unregistered reader'' failure path. A real application must register
the reader families it intends to use before any \texttt{.xnb} load. CNA supplies one explicit,
idempotent umbrella call for every implemented built-in family (but deliberately not arbitrary
closed generic collection readers):

\begin{lstlisting}
#include "CNA/Internal/Xnb/XnbBuiltInReaders.hpp"

CNA::Internal::Xnb::RegisterAllBuiltInXnbReaders(); // once, before any .xnb load

ContentTypeReaderManager::AddTypeCreator(
    "MyGame.Content.GameLevelDataReader",
    [] { return std::make_unique<GameLevelDataReader>(); });

ContentManager content(nullptr, "Content");
GameLevelData level = content.Load<GameLevelData>("level1");
\end{lstlisting}

The custom-reader integration test really writes that minimal \texttt{level1.xnb}, registers the
factory, and loads it through \texttt{ContentManager}; omitting the registration yields the
named \texttt{ContentLoadException}. This is the route a port of
\texttt{prime31/Nez}'s independent \texttt{BitmapFontReader} must take. CNA deliberately has no
equivalent to .NET's reflection-generated \texttt{ReflectiveReader<T>}: content built around one
must be rebuilt with an explicit writer/reader pair, then supplied through a matching C\texttt{++}
factory.

This global registry has deliberately different replacement and lifetime rules from the
per-manager tables. \texttt{AddTypeCreator()} keeps the \emph{first} factory for an exact name
and silently ignores later registrations; \texttt{CreateReader()} calls the retained factory for
one fresh reader per \texttt{.xnb} file, owned by that file's \texttt{ContentReader}.
\texttt{ClearTypeCreators()} removes every family for every manager and is consequently a test
isolation tool, not a runtime reload API. Registration and clearing have no synchronization in
the current C\texttt{++} map, so finish them before concurrent asset loads. There is also no
empty-name or empty-factory validation: the manifest test deliberately stores a null factory and
then observes \texttt{IsRegistered()} report true; attempting to create through that entry would
invoke an empty \texttt{std::function} rather than yield a normal missing-reader diagnostic.

\texttt{RegisterTypeReader<T>}, by contrast, changes only one \texttt{ContentManager}'s
loose-file/\texttt{.cnj} reader map. It owns the supplied \texttt{unique\_ptr} as a
\texttt{shared\_ptr} and silently overwrites an existing reader for the same C\texttt{++} type;
it accepts a null pointer, which later reaches the map's explicit ``reader is null'' load error.
It neither changes the global \texttt{.xnb} registry nor evicts an asset already cached under
that type/name, so call \texttt{Unload()} before expecting a replacement reader to affect a
previously loaded generic asset. No located test directly registers and replaces a loose reader.

\texttt{RegisterCnjLoader<T>(typeName, factory)} is safer but has its own narrower contract. On
its first successful call for a type with no reader, it creates the one generic \texttt{.cnj}
reader; later distinct names select separate factories inside that same reader. Empty names,
empty factories, a pre-existing non-generic reader, and an exact duplicate name all fail fast,
and the focused tests cover those cases plus recursive texture loading. Do not subsequently
replace that generated reader with \texttt{RegisterTypeReader<T>}: current CNA leaves the named
factory table behind, so already registered and even later accepted names become unreachable
through the replacement reader. This is an untested implementation boundary, not a supported
way to layer two loose-reader mechanisms. Copying a \texttt{ContentManager} shares its current
loose-reader pointers but copies its named \texttt{.cnj} factories, as the later
\texttt{Game::Content} section explains.

\section{One object-graph session}

\subsection{\cnaclass{ContentReader} is a one-shot XNB object-graph session, not a reusable file reader}
\label{sec:content-reader-contract}

Every real XNB load eventually creates a \cnaclass{ContentReader} over the body beginning
immediately after the ten-byte container header. Its first operation, \texttt{ReadAsset<T>()},
parses the type-reader table at the \emph{current} stream position, reads the shared-resource
count, deserializes one root object, then reads every shared object and finally invokes its
queued fixups. The constructor comment currently says that both the header \emph{and} table have
already been consumed, which contradicts the implementation and the manager's live call site:
only the header has been consumed. Calling \texttt{ReadAsset<T>()} a second time does not rewind
or reuse the prior table; it treats whatever bytes follow the first asset as another table.
Use one new reader for one body and one top-level asset.

Although FNA makes this a sealed class with an internal constructor, CNA's C\texttt{++} class,
constructor, and the initialization/read methods are public and it is not final. That wider
surface is not a promise that the object is independently long-lived. The reader stores raw,
non-owning pointers to both stream and manager. Its inherited \cnaclass{BinaryReader} uses the
default \texttt{leaveOpen = false}: destroying or closing the reader closes the supplied stream
but never deletes it. The manager-created stack reader is therefore safe only because its stack
\cnaclass{MemoryStream} outlives it. A standalone caller must keep both raw referents alive,
must not retain a \texttt{ContentReader\&} beyond its callback, and must account for the stream
being closed on reader destruction. Neither the focused tests nor the public API prevent a
dangling manager/stream or concurrent use of the mutable reader state.

Initialization creates a fresh reader instance for every serialized table entry through the
process-global factory registry, checks CNA's extra strict serialized-version rule, and only then
calls \texttt{Initialize()} on every instance. FNA reads but does not enforce reader versions.
The two-pass construction is useful: a reader may assume its peers have all been created when its
initializer runs. There is a C\texttt{++}-specific lifetime caveat, however. CNA passes that
initializer a local \cnaclass{ContentTypeReaderManager} which is destroyed as soon as
initialization completes; a custom reader must use the reference synchronously and must not store
it. FNA retains its manager as a field of the ContentReader for the reader session. The fresh CNA
reader instances themselves are also destroyed when the session ends, so an asset must not retain
a pointer/reference to its decoder.

Object dispatch is 1-based: zero is a null/default object, while positive index $n$ selects table
entry $n-1$. For a default-constructible C\texttt{++} result, zero returns \texttt{T\{\}}; with
an existing instance it returns that instance instead. A non-default-constructible \texttt{T}
with no existing value has no C\#-style \texttt{default(T)} representation and throws
\cnaclass{ContentLoadException}. A bad reader/result pairing instead exposes
\texttt{std::bad\_any\_cast}; CNA does not normalize every object-graph failure. Direct
\texttt{ReadObject} and \texttt{ReadRawObject} overloads can bypass the 1-based dispatch, but
CNA lacks FNA's no-argument \texttt{ReadRawObject<T>()} lookup overload. Root success is not a
transaction: a later shared-resource parse or fixup can throw after the reader has already made
arbitrary reader-side changes.

Shared resources deliberately have a second phase. A positive reference index queues a typed
fixup; index zero queues nothing. After the root returns, CNA reads \emph{all} shared values in
file order, then runs each resource's fixups in file/index order. This correctly supports forward
and cyclic graph references, but it also means a callback observes a fully-read shared section,
not an incrementally available one. A positive index above the declared count produces a content
load error. A negative index instead takes the same silent no-fixup branch as zero, because the
code tests only \texttt{index > 0}. A wrong typed shared value reaches
\texttt{std::bad\_any\_cast}; current FNA turns that mismatch into its content-load error. The
focused tests prove one healthy fixup and zero, but not a negative index, bad type, callback
throw, forward/cyclic graph, or partial-side-effect case.

Disposal tracking is narrower than the constructor documentation promises. CNA records a result
only when its \emph{static} C\texttt{++} type is \texttt{std::shared\_ptr<U>} for a
\texttt{U} derived from \cnaclass{IDisposable}, and only when an explicit callback was supplied.
The normal \cnaclass{ContentManager} XNB path supplies no callback; despite the header's statement
about a manager fallback, the live implementation has no such call. Type-erased shared-resource
reads deliberately skip tracking altogether. FNA instead performs a runtime IDisposable test and
falls back to its owning ContentManager. No located CNA test exercises either a recording callback,
the documented fallback, or a disposable shared resource, so do not use the reader as a general
content-disposal coordinator.

\texttt{ReadExternalReference<T>()} has two separate CNA restrictions. An empty serialized
reference returns \texttt{std::nullopt}; a nonempty one needs a live manager, resolves relative to
the containing logical asset name, and then enters the ordinary \texttt{Load<T>()} candidate and
cache path. Relative upward escape is rejected, but an absolute reference still escapes through
the manager's unrestricted absolute-path route described in \S8.4. The declaration looks generic,
but the implementation is explicitly instantiated only for \cnaclass{Texture2D} and
\cnaclass{TextureCube}; another \texttt{T} compiles against the declaration but has no linked
definition. FNA's method is genuinely generic and returns \texttt{default(T)}; CNA's
\texttt{optional} return is the necessary value-type distinction. Existing tests cover only the
Texture2D sibling, empty, relative-escape, and null-manager paths.

Finally, \cnaclass{XnbReadLimits} is not a universal admission policy on the current
\texttt{develop} source. Type-reader and shared-resource counts are checked; collection and
decoded-pixel limits work when a particular built-in reader calls the corresponding helper.

Three declared limits have gaps: \texttt{maxStringBytes} and
\texttt{maxObjectNestingDepth} have no live consumer, while \texttt{maxFileSize} is consulted
only on the LZX decompression path. An uncompressed XNB is first read wholesale by
\cnaclass{ContentManager}, and deeply nested object dispatch or a nested generic type-reader name
has no reader-level depth guard. Treat the defaults as partial defensive checks, not proof that
every untrusted XNB has a file, string, allocation, or stack-depth bound. No located focused test
proves those three advertised limits on this branch.

\subsection{\cnaclass{ContentTypeReader} is a value-erased adapter, not a runtime type system}
\label{sec:content-type-reader-contract}

CNA keeps the familiar template name \texttt{ContentTypeReader<T>} for a game author's typed
base class. Its separate, \texttt{CNAEXT} \cnaclass{ContentTypeReaderBase} exists because C\texttt{++}
cannot give both the non-template and template forms the same unqualified name as C\# can. A
concrete reader implements the protected typed \texttt{Read(ContentReader\&, optional<T>)}; the
base translates that call to public \texttt{ReadUntyped(ContentReader\&, any)}. The stored target
type is only a caller-supplied canonical XNA \emph{name} string. CNA has neither reflection nor a
runtime \texttt{Type} identity, and it does not verify that the string, the registered table name,
the template \texttt{T}, and the value actually returned by the reader agree. The eventual
\texttt{any\_cast} is the real check, and a mismatch exposes \texttt{std::bad\_any\_cast}, not a
uniform content-load exception.

For a copy-constructible \texttt{T}, the erased payload is a literal \texttt{T}. For a move-only
\texttt{T}, it is instead \texttt{shared\_ptr<T>}, because \texttt{std::any} itself requires a
copy-constructible held type. The normal indexed \texttt{InnerReadObject<T>()} path contains the
matching compile-time unwrap and therefore loads bare move-only roots correctly. The public direct
\texttt{ReadObject<T>(reader)} and \texttt{ReadRawObject<T>(reader)} overloads do \emph{not}:
they always cast the returned \texttt{any} to bare \texttt{T}. For a move-only value they therefore
reach \texttt{bad\_any\_cast}; their overloads taking an existing value try to put that bare value
in \texttt{any} and cannot be instantiated. This affects both \cnaclass{SoundEffect} and
\cnaclass{TextureCube}, not only the former. Use ordinary indexed dispatch for those types; the
direct-reader overloads are currently reliable only for a copy-constructible target (or a reader
whose declared \texttt{T} is already a copyable handle such as \texttt{shared\_ptr<Texture3D>}).

An existing instance is likewise a transported value, not an observable in-place update of the
caller's object. CNA moves it into \texttt{optional<T>}, lets the reader return a replacement
value, then returns that value to the caller. In particular, no production CNA call site consults
\texttt{getCanDeserializeIntoExistingObjectProperty()}: normal \texttt{ContentManager} loading
starts without one and caches its completed result. Only \texttt{ListReader} and
\texttt{DictionaryReader} advertise true, but a direct caller still receives a new C\texttt{++}
container value. The list result appends decoded items after the supplied values; the dictionary
result clears its supplied entries first. \texttt{ArrayReader} advertises the default false yet
accepts an explicitly supplied vector and resizes it to the file count, a deliberate safety
departure from FNA's unchecked array indexing. No collection test proves the public
\texttt{ContentReader} overloads, a real manager reload, or the distinction between a returned
replacement and mutation of a caller-owned container.

Collection genericity is also deliberately finite. FNA's per-file manager uses reflection and
\texttt{TargetType} to find the element reader from the current table during \texttt{Initialize()}.
CNA instead requires each closed \texttt{ArrayReader}, \texttt{ListReader},
\texttt{DictionaryReader}, or \texttt{NullableReader} registration to carry canonical names for
its element readers. At every collection read it creates those readers afresh from the global
factory map; they are neither initialized from the current file nor checked against a serialized
version. This is sound only for the currently stateless built-ins, as the source comment says.
The umbrella registration deliberately excludes arbitrary closed generics; only SpriteFont's three
known list shapes are pre-registered. A custom collection must register its exact outer
combination and every value-type element reader before loading, rather than expecting a
reflection-like fallback.

For \texttt{shared\_ptr} elements the collection consumes one indexed object dispatch per item,
allowing the file table to select a concrete subtype; non-\texttt{shared\_ptr} elements use the
manually named direct reader. All three counted collections check the configured element limit
before allocation or reserve, and the focused tests cover unsigned overflow, negative counts, and
one above-limit count. They do not cover reference-shaped elements, a stateful custom element
reader, factory exceptions/null returns, or duplicate dictionary keys. The latter diverges from
FNA particularly: FNA uses \texttt{Dictionary.Add} and throws for a duplicate, while CNA uses
\texttt{unordered\_map::emplace} and silently retains the first value.

The built-in marker reader for the general EffectReader is the other intentional branch of this
framework. It consumes no input and always gives a named content-load error for compiled platform
shader bytecode. This is a precise scope boundary, not a generic reader fallback. Tests prove its
registration and direct throw, plus healthy custom registration and selected collection behavior;
none exercises incorrect \texttt{any} boxing, the move-only direct-overload trap, target-name
disagreement, initializer lifetime, or the global registry's null-returning factory.

\section{Built-in reader families}

\subsection{The scalar and geometry readers have literal layouts, but no enum fallback}
\label{sec:xnb-scalar-math-reader-contract}

Two intentionally small reader families form the value-type baseline: thirteen primitives and
thirteen framework/math values. Their internal C\texttt{++} class names are not part of the XNB
protocol; \texttt{RegisterPrimitiveXnbReaders()} and \texttt{RegisterMathXnbReaders()} install
the real \texttt{Microsoft.Xna.Framework.Content.*Reader} keys. Calling the umbrella
\texttt{RegisterAllBuiltInXnbReaders()} includes both families, but none is registered implicitly
by constructing a \cnaclass{ContentManager}. Like every reader that does not override the common
method, all twenty-six accept only serialized version zero under CNA's stricter
\texttt{SupportsVersion()} policy. FNA reads and discards those table version integers, so a
nonzero version on an otherwise familiar primitive is a deliberate CNA compatibility rejection.

The primitive layouts are exactly the underlying \cnaclass{BinaryReader} operations, rather than
host-memory copies: Boolean and signed/unsigned bytes consume one byte (any nonzero Boolean byte
is true); 16-, 32-, and 64-bit integers are little-endian; \texttt{Single} and
\texttt{Double} are little-endian IEEE-754 bit patterns. A \texttt{Char} is different from a
two-byte C\texttt{++} \texttt{char16\_t}: it is one UTF-8-decoded UTF-16 code unit and may consume
one to three input bytes. A four-byte code point is rejected because it requires a surrogate pair.
The text reader takes a 7-bit byte length followed by those raw bytes into \texttt{std::string};
it does not separately validate UTF-8 or apply \texttt{maxStringBytes}. Thus an XNB field called
\texttt{String} is a length-delimited byte string at this boundary, with the same remaining-stream
preflight and allocation caveat described for type-reader names above. Truncation and malformed
character/length encodings retain the underlying binary-reader exceptions, not a reader-specific
\cnaclass{ContentLoadException}.

The geometric layouts are equally literal: \texttt{Vector2}, \texttt{Vector3}, and
\texttt{Vector4} are two, three, and four Singles in component order; \texttt{Quaternion} is
\texttt{x,y,z,w}; and \texttt{Matrix} is sixteen Singles in the source field order
\texttt{M11} through \texttt{M44}. \cnaclass{Color} is four consecutive \texttt{r,g,b,a} bytes,
not its packed integer representation. \cnaclass{Plane} is Vector3 normal then \texttt{D};
\cnaclass{Point} is two Int32 values; and \cnaclass{Rectangle} is left, top, width, height as
four Int32 values. A bounding box is min Vector3 followed by max Vector3, a sphere is center
Vector3 followed by radius, a frustum is reconstructed from one Matrix, and a ray is position
Vector3 followed by direction Vector3. These readers merely construct their value after consuming
the fields: they do not add finite-value, normal-length, radius, min/max, rectangle-size, or
matrix-invertibility validation. Any invariant a framework value enforces is therefore separate
from XNB admission.

There is no corresponding generic \texttt{EnumReader<T>} in current CNA. The source tree has no
enum reader class, registration, or specialized factory key, and the umbrella registration cannot
supply one. FNA's own reader manager explicitly instantiates generic enum readers (including
\texttt{SpriteEffects} and \texttt{Blend}) for its reflection/AOT machinery. CNA's type-name
normalizer can parse a generic EnumReader table name, but the later global factory lookup reports
it as unregistered. A compiled asset whose root or nested object relies on a generic enum reader
therefore cannot load without an application-provided exact canonical factory; a C\texttt{++}
enum's underlying type alone is not a fallback policy.

The focused primitive/math test writes values through the same local \cnaclass{BinaryWriter} and
calls \texttt{ReadUntyped()} directly, without an XNB header, table, dispatch index, fixture from
another producer, truncation case, or nonzero-version case. It checks every primitive reader, but
only Vector2, Color, Plane, Point, Rectangle, BoundingBox, BoundingSphere, and Ray among the
math layouts; Vector3, Vector4, Matrix, Quaternion, and BoundingFrustum have registration-only
coverage there. The two real MonoGame SpriteFont end-to-end fixtures additionally exercise
Char, Rectangle, and Vector3 through nested lists, but no located fixture proves the remaining
layouts, Unicode edge cases, Boolean false/noncanonical true, malformed UTF-8, string bounds, or
the enum gap. Treat this family as a transparent binary compatibility layer for the stated
layouts, not as a validating schema or complete XNB value-type surface.

\subsection{Decimal and calendar readers preserve ticks more reliably than all metadata}
\label{sec:xnb-decimal-datetime-reader-contract}

\texttt{RegisterDecimalDateTimeXnbReaders()} supplies \texttt{TimeSpanReader} and
\texttt{DateTimeReader} on every supported C\texttt{++} compiler, but supplies
\texttt{DecimalReader} only when \texttt{\_MSC\_VER} is absent. That is a real platform
compatibility boundary, not an optional fast path: sharp-runtime's \cnaclass{Decimal} uses
\texttt{unsigned \_\_int128}, so its header, reader class, and canonical factory are deliberately
compiled out for the MSVC family. An XNB that names DecimalReader consequently reaches the ordinary
unregistered-type-reader error on that build, while the TimeSpan and DateTime names remain
available. As with the preceding value readers, the three inherit exact serialized version zero
and require explicit registration (or the umbrella helper) before a manager load.

Decimal consumes sixteen little-endian bytes in .NET wire order: signed Int32 \texttt{lo},
\texttt{mid}, \texttt{hi}, then \texttt{flags}. The first three are reassembled as an unsigned
96-bit mantissa; flag bit 31 supplies the sign, and bits 16--23 supply the decimal scale. The
sharp-runtime \cnaclass{Decimal} constructor correctly rejects a scale above 28, so that failure
escapes as its argument-range exception rather than as \cnaclass{ContentLoadException}. It does
\emph{not}, however, validate the remaining reserved flag bits. CNA silently discards them, and
also canonicalizes a signed zero to an unsigned zero. A small reference run through Mono's
\texttt{BinaryReader.ReadDecimal()} rejects a flag word with reserved bit 0 set as
\texttt{IOException} and preserves a negative-zero flag; current FNA delegates to that runtime
operation. CNA therefore accepts some corrupt Decimal payloads that the reference rejects and
does not reproduce every observable Decimal bit pattern, even though its ordinary positive
mantissa/scale layout is correct.

TimeSpan is the uncomplicated case: it is one signed little-endian Int64 count of 100-nanosecond
ticks, passed directly to \cnaclass{TimeSpan}; every 64-bit signed value is representable. DateTime
instead reads an unsigned 64-bit word, treats the low 62 bits as ticks and the high two as
\texttt{DateTimeKind}, then constructs sharp-runtime \cnaclass{DateTime} from ticks alone. The
constructor still rejects a tick count below zero or above its year-9999 maximum, but the kind is
lost because sharp-runtime does not store it or offer local/UTC conversion APIs. This loses valid
Unspecified/UTC/Local distinctions. It also changes malformed input admission: the code masks and
discards kind value 3, whereas FNA passes it to \texttt{new DateTime(ticks, kind)}; the same Mono
reference run accepts 0, 1, and 2 but throws \texttt{ArgumentException} for 3. CNA accepts that
otherwise valid-tick word as an unspecified value.

The focused test shares the direct-writer/direct-\texttt{ReadUntyped()} shape of the primitive
test: one positive Decimal with only \texttt{lo} and scale checked, one positive TimeSpan, and
one UTC-bit DateTime. It intentionally conditionalizes Decimal registration on the local compiler,
but supplies no MSVC build result, external XNB fixture, full 96-bit/sign/negative-zero Decimal,
reserved flag, invalid scale, negative/extreme TimeSpan, Local/Unspecified/3 DateTime-kind, invalid
tick, truncation, or version test. These readers make common values usable, but they are not a
cross-platform, bit-for-bit fidelity guarantee for all XNA temporal and Decimal payloads.

\subsection{CurveReader preserves FNA's field order but not a curve-input policy}
\label{sec:xnb-curve-reader-contract}

The separately registered \texttt{Microsoft.Xna.Framework.Content.CurveReader} is the binary
counterpart of Chapter~\ref{ch:math-core-types}'s \texttt{Curve} discussion, not the JSON
\texttt{.cnj} reader illustrated there. Its target name is
\texttt{Microsoft.Xna.Framework.Curve}; it joins the built-in set only after an explicit
\texttt{RegisterCurveXnbReader()} call (or the umbrella registration call), and it inherits the
usual exact serialized-version-zero requirement. Current FNA's internal \texttt{CurveReader}
uses the same field sequence:

\begin{center}
\begin{tabular}{r p{0.76\linewidth}}
  1 & \texttt{Int32 PreLoop} (\texttt{CurveLoopType}) \\
  2 & \texttt{Int32 PostLoop} (\texttt{CurveLoopType}) \\
  3 & signed \texttt{Int32 keyCount} \\
  4 & for each key: \texttt{Single position}, \texttt{value}, \texttt{tangentIn},
      \texttt{tangentOut}, then \texttt{Int32 continuity} (\texttt{CurveContinuity}).
\end{tabular}
\end{center}

The word ``tangent'' in that layout is important: the four key numbers include two already
calculated floating-point tangents. They are not \texttt{CurveTangent} enum values, and the
reader does not invoke \texttt{ComputeTangents()}. The three integer enum sites are merely
\texttt{static\_cast}s. Valid loop codes are Constant=0, Cycle=1, CycleOffset=2,
Oscillate=3, and Linear=4; valid continuity codes are Smooth=0 and Step=1. No range check
occurs before either value enters the returned curve. An invalid continuity is subsequently
treated as Smooth, because evaluation tests only equality with Step. An invalid pre- or post-loop
value reaches the switch with no matching case and falls through to ordinary
\texttt{GetCurvePosition(position)} rather than producing a content-load error. Its result then
depends on the positions and tangents present in that particular curve. These are malformed-file
admission details, not alternate documented enum meanings.

Nor does this reader validate the five floating-point fields or impose key semantics beyond what
\cnaclass{CurveKeyCollection} happens to do. NaN/infinite numbers, duplicate positions, and
non-monotone serialized order are accepted. \texttt{Add()} insertion-sorts only when a new
position compares strictly less than an old one, so equal positions retain input order and NaN
positions append. Subsequent curve evaluation can therefore encounter a zero position span or
non-finite arithmetic; this reader neither rejects nor repairs it. This is deliberately unlike
the textual \texttt{.cnj} examples' named enum vocabulary: a binary cast is not validation.

The declared key count is the sharper input boundary. CNA reads it as signed \texttt{Int32} and
immediately uses \texttt{for (i = 0; i < keyCount; ++i)}. It does \emph{not} call
the common collection-count helper, despite that helper rejecting a
negative count and enforcing the configured ten-million-element default before collection
growth. Thus a negative count succeeds as an empty key sequence after the two loop fields, while
a very large positive count repeatedly appends keys until the body stream fails or storage is
exhausted. It does not reserve the full count in advance, but it has no reader-level count ceiling
either. A current hardening commit that notes this omission exists only on the unmerged
\texttt{feature/audit} branch; it is not behavior supplied by this book's current
\texttt{develop} authority.

There is also an easy-to-miss existing-instance asymmetry. CurveReader does not advertise
\texttt{CanDeserializeIntoExistingObject}; the normal \cnaclass{ContentReader} object dispatch
therefore always passes no existing curve. A direct type-erased caller can nevertheless supply a
\cnaclass{Curve}. CNA takes that value copy, overwrites its two loop fields, and appends newly
read keys without clearing its old ones; the returned C\texttt{++} value is separate from the
caller's original. FNA's same append logic operates on its reference-type instance. This corner
does not affect a normal manager load, but it means direct reader reuse is neither a replacement
operation nor a caller-object mutation contract in CNA.

The focused CurveReader test writes two ordinary, locally generated keys through CNA's own
\texttt{BinaryWriter}, calls \texttt{ReadUntyped()} directly, and checks registration, both loop
fields, positions, values, selected tangents, and the two valid continuities. It has no XNB
header/table dispatch, external producer fixture, manager load, existing-instance, zero/negative/
large count, truncation, nonzero version, invalid enum, duplicate/NaN/infinite key, or evaluation
test. The checked-in XNB fixture inventory contains no Curve asset; the related real fixture is
only the \texttt{.cnj} Curve covered in Chapter~\ref{ch:math-core-types}. Treat the binary
reader as a faithful ordinary-layout decoder, not as a bounded or semantic curve validator.

\subsection{Texture2DReader validates decoded Color levels but not texture topology}
\label{sec:xnb-texture2d-reader-contract}

\texttt{Microsoft.Xna.Framework.Content.Texture2DReader} is separately registered (or included
by the umbrella call), targets \cnaclass{Texture2D}, and inherits the common exact
serialized-version-zero rule. Its payload begins with four signed 32-bit words:
\texttt{surfaceFormat}, \texttt{width}, \texttt{height}, and \texttt{levelCount}; every declared
level then has a signed \texttt{byteCount} followed by that many bytes. The first word depends on
the \emph{XNB container} version rather than the reader-table version. Version 5 and later are
cast directly to \cnaclass{SurfaceFormat}. Version 4 accepts only the FNA legacy ordinals 1
(\texttt{ColorBgraEXT}), 28 (DXT1), 30 (DXT3), and 32 (DXT5). Unknown values fail as a
\cnaclass{ContentLoadException}; no general enum validation takes place before that later
format decision.

The practical format surface is deliberately smaller than that decoding step makes it appear.
CNA accepts raw \texttt{Color}, and accepts DXT1, DXT3, and DXT5 only by software-decompressing
them to RGBA \texttt{Color} before upload. Every other current \cnaclass{SurfaceFormat} is
rejected, including \texttt{ColorBgraEXT}; consequently the otherwise recognized version-4
legacy ordinal 1 immediately reaches CNA's ``only Color/Dxt1/Dxt3/Dxt5'' error. Current FNA
retains a native DXT format when the active FNA3D device advertises DXT/S3TC support and converts
only when it cannot; CNA always consumes the larger uncompressed representation. That is a
performance and memory difference rather than a reason to expect different ordinary pixels.
It is nevertheless a real compatibility limit for the many non-Color XNA surface formats.

There is a second platform-specific difference. FNA rewrites Color/ColorBgra and the three DXT
encodings when the XNB platform byte is \texttt{'x'} (Xbox 360) before conversion/upload. CNA's
container accepts that platform code, but this reader never performs the corresponding
\texttt{X360TexUtil} swaps. A supported-looking Xbox texture can therefore decode with wrong
channel/block layout even when its nominal format is one of the four accepted forms. The fixture
tree has no version-4 or Xbox Texture2D XNB with which to establish an interoperable result.

The existing byte defenses are useful but deliberately local. Width and height must each be
positive, then the reader calls \texttt{CheckDecodedByteSize(width * height * 4)} before creating
the texture; with the default \cnaclass{XnbReadLimits}, that nominal base-level RGBA size is
capped at 256 MiB. For an uncompressed Color level, it reads exactly the declared byte count and
then requires exactly \texttt{levelWidth * levelHeight * 4} bytes before constructing
\cnaclass{Color} objects. A negative byte count becomes \cnaclass{ContentLoadException}; a short
stream remains \texttt{EndOfStreamException}. DXT's helper checks that there are \emph{at least}
the required 4-by-4 blocks (eight bytes per DXT1 block and sixteen per DXT3/5 block), then emits
exactly one RGBA image. It does not reject surplus compressed bytes in a declared level, and its
too-short-data failure is \texttt{std::out\_of\_range}, not the reader's content-load error.
Thus this is not a uniform malformed-texture exception policy.

Nor is the initial product a complete allocation or device admission policy. It has no per-axis
device maximum, so a shape such as 500000 by 1 has a small nominal byte count yet can be invalid
for the selected GPU. More subtly, the current source uses raw signed \texttt{int64\_t}
multiplication: two legal positive Int32 dimensions can leave \texttt{width * height} representable
but overflow on the following \texttt{* 4} before \texttt{CheckDecodedByteSize()} observes a
value. That is C\texttt{++} signed-overflow undefined behavior, not a clean limit rejection.
The unmerged \texttt{feature/audit} branch contains a checked-multiply and a device-axis preflight,
but neither commit is an ancestor of this book's \texttt{develop} authority. It must not be
treated as shipped protection.

Mip topology has the same distinction between a normal asset and hostile input. Construction uses
only \texttt{levelCount > 1}: a positive multi-level declaration creates the framework's full
halving chain, rather than a chain exactly as long as the file says. Zero or negative counts make
a one-level texture and run no per-level reads, so this reader accepts an uninitialized texture
payload instead of rejecting an invalid count. A positive count above the physical chain has no
ceiling. The reader shifts \texttt{width >> level} and \texttt{height >> level}, while
\cnaclass{Texture2D::SetData()} checks only that the level is nonnegative; its CPU mip store grows
for arbitrary higher levels and the renderer is asked to update them. Large levels also make the
shift itself undefined. An unmerged audit change adds an upper mip ceiling, but current
\texttt{develop} does not have it. The base-level decoded-size check is therefore not a bound on
the number, topology, or total storage of declared mip payloads.

A GPU device is mandatory for ordinary manager loading. With a non-null \cnaclass{ContentManager}
that lacks one, \texttt{getGraphicsDeviceInternal()} throws the manager's own diagnostic before
the reader's null-device branch. The latter is relevant only to a manually constructed
\cnaclass{ContentReader} with no manager. Direct existing-instance use has another non-obvious
shape: normal object dispatch passes none, but a direct boxed \cnaclass{Texture2D} is moved into
the reader and reused without checking its dimensions, format, device, or mip count against the
file. File dimensions determine the byte check, whereas \texttt{SetData()} uses the retained
texture's dimensions. The returned C\texttt{++} handle is the moved/reused value, not a guarantee
that a caller-owned object was overwritten as a transactional replacement. FNA likewise reuses a
non-null reference, but has managed reference identity rather than this value-handoff boundary.

The focused unit test proves registration, unsupported \texttt{Bgr565}, one bad Color byte count,
negative width, and two-axis huge dimensions. It does not prove a successful direct Color or DXT
decode, legacy/version-4 mapping, Xbox swapping, zero height/width, negative/zero/excessive mip
counts, one-axis device limits, overflow arithmetic, truncated/negative/surplus level data,
missing device, existing-instance mismatch, or reader version. End-to-end coverage is better but
narrow: a real MonoGame 1-by-1 white Color XNB verifies its exact pixel; a real LZX-compressed
64-by-64 Color explosion is checked only for non-uniform output; and the real SpriteFont fixture
exercises a nested 128-by-128 DXT3 atlas without independently checking its pixels. There is no
checked-in Texture2D DXT1/DXT5, version-4, Xbox, or non-Color fixture. Treat the reader as a
well-tested ordinary Color path with some real compressed integration, not as complete
Texture2D/XNB format or adversarial-GPU validation.

\subsection{Texture3DReader now slices DXT volumes, but does not admit a physical volume}
\label{sec:xnb-texture3d-reader-contract}

The separately registered \texttt{Microsoft.Xna.Framework.Content.Texture3DReader} is unusual in
the C\texttt{++} family because its target is \texttt{std::shared\_ptr<Texture3D>}, not a
bare resource value. That makes it copyable enough for \texttt{std::any}, while the resource still
owns a non-copyable renderer handle. It inherits exact serialized version zero. CNA and current
FNA consume the same literal payload: Int32 \texttt{surfaceFormat}, \texttt{width},
\texttt{height}, \texttt{depth}, and \texttt{levelCount}, followed by one Int32 \texttt{byteCount}
and payload for each level. Unlike Texture2DReader, there is no pre-version-5 format remapping;
both readers cast the format word directly and halve all three dimensions after each level.
FNA's mip-chain constructor nevertheless counts only width and height, not depth, and CNA's
\cnaclass{Texture3D} deliberately follows that same rule.

The ordinary format boundary is narrower than FNA's. FNA forwards the file's raw level bytes and
nominal \cnaclass{SurfaceFormat} to its Texture3D implementation. CNA accepts only raw
\texttt{Color} or DXT1/DXT3/DXT5, converting the latter to RGBA Color before its own upload;
every other format, including a recognized but unsupported enum value, fails as
\cnaclass{ContentLoadException}. There is no Texture2D-style legacy mapping or Xbox byte swap in
either FNA's or CNA's Texture3D reader. CNA's unconditional DXT conversion has the same portability
versus native-memory/performance trade-off as its 2D path, but its much smaller accepted format
set remains a compatibility limit.

One important current-develop hardening fix is real. A compressed volume is a sequence of
independently compressed two-dimensional slices, not one DXT image with a very tall height. The
reader divides a level's payload by its current depth, feeds each slice separately to
\texttt{DxtUtil::DecompressDxt1/3/5}, and concatenates the RGBA slices. The prior whole-volume
interpretation discarded later slices; it was found by the broad XNB hardening sweep and is fixed
on this branch. The boundary is still permissive: it does not require the byte count to be evenly
divisible by depth or exactly equal to the required block total. Each helper rejects a slice that
is too short, but accepts extra full-slice data, and a remainder after the division has already
been consumed from the stream but is never decompressed. DXT underflow is still
\texttt{std::out\_of\_range}, while a negative declared count is \cnaclass{ContentLoadException}
and a short physical stream is \texttt{EndOfStreamException}.

Raw Color levels receive a useful exact check after reading: the bytes must equal the current
\texttt{width * height * depth * 4} before the reader builds real \cnaclass{Color} objects. Base
dimensions must also be individually positive and are presented to
\texttt{CheckDecodedByteSize()} before texture creation. As in the 2D reader, that is not a
complete arithmetic guard in current \texttt{develop}. The initial raw signed
\texttt{int64\_t width * height * depth * 4} expression can overflow before the limit helper
observes it, and the later \texttt{int32\_t voxelCount = width * height * depth} is an independent
overflow site. The checked-multiply and voxel-range corrections exist in unmerged
\texttt{feature/audit}, not in the authority documented here. A valid-looking total byte size
also supplies no per-axis device maximum or reader-level aggregate mip-data budget.

Mip admission has a separate hole. Construction receives only \texttt{levelCount > 1}, so a
positive multi-level file allocates the framework's complete width/height chain rather than a
chain exactly as long as its declaration; zero or negative counts create a one-level texture and
perform no data read. An excessive positive count is not rejected. After the genuine chain ends,
the reader continues with 1-by-1-by-1 payloads and calls \texttt{SetData()} for arbitrary level
numbers. Texture3D checks only a negative level and relative box ordering; it neither checks that
the level exists nor that the box fits the actual stored mip dimensions before forwarding to its
renderer. Thus the base decoded-size cap neither validates a serialized mip topology nor protects
all renderer update paths.

A \cnaclass{GraphicsDevice} is required only when the reader must create a texture. A non-null
manager without one throws the manager's own diagnostic; a manually built ContentReader with no
manager reaches the reader's error. But a direct caller may supply a non-null
\texttt{shared\_ptr<Texture3D>}; normal dispatch never does. That pointer is copied and the same
referenced resource is mutated in place, without a format, device, dimensions, or mip-chain match
against the file; it also bypasses the reader's new-texture device lookup. This is close to FNA's
reference-instance behavior, but notably unlike CNA Texture2D's transported value. Renderer
support is not checked at either point: the default 3D renderer factory can return null, after
which Texture3D's upload and readback paths simply skip renderer work. Loading an XNB can therefore
return a nominal Texture3D whose supplied voxels were not represented by a real renderer resource
on a 2D-only implementation.

Evidence for the ordinary wire is narrow. The joint 3D/cube unit file proves registration,
rejects one unsupported \texttt{Bgr565} word, and directly feeds a locally written 2-by-2-by-1
Color body to \texttt{ReadUntyped()}, then reads four exact voxels back. There is no external
Texture3D XNB anywhere in the fixture tree, no header/table/ContentManager XNB test, and no
direct DXT, multi-slice, mip, truncation, dimensions, byte-count, version, device, existing-pointer,
or renderer-capability case. The whole-container fuzzer has real Model, Texture2D, and SoundEffect
seeds only. The separate CNJ test does exercise a self-contained 2-by-2-by-2 volume, but it is a
different loader and cannot establish XNB compatibility. Treat the current reader as a faithful
ordinary field-order implementation with a repaired slice algorithm, not as externally-fixtured
or universally renderer-safe volume-texture support.

\subsection{TextureCubeReader loads a real DXT1 cube but still trusts raw Color byte counts}
\label{sec:xnb-texturecube-reader-contract}

The explicitly registered \texttt{Microsoft.Xna.Framework.Content.TextureCubeReader} has FNA's
simple wire order: Int32 \texttt{surfaceFormat}, square Int32 \texttt{size}, and Int32
\texttt{levels}; then, in this exact nested order, six faces and every level of each face, an
Int32 byte count followed by its bytes. Face values 0 through 5 map directly to
\texttt{PositiveX}, \texttt{NegativeX}, \texttt{PositiveY}, \texttt{NegativeY},
\texttt{PositiveZ}, and \texttt{NegativeZ}. The file does not serialize a per-face dimension:
CNA starts every face at the base size and halves it after each level. Neither current FNA nor CNA
has a Texture2D-style legacy version mapping or Xbox byte-swapping branch here. CNA again applies
the common exact serialized-version-zero rule around FNA's ordinary payload decoding.

FNA forwards each declared blob and its nominal \cnaclass{SurfaceFormat} to TextureCube. CNA
accepts Color plus DXT1/DXT3/DXT5 only, always software-decompressing DXT to Color, and rejects
all other formats. This is the same format-surface reduction as the other texture readers. For
compressed levels, the current DXT helper verifies that the data contains at least the required
4-by-4 block sequence (including a full minimum block for 2-by-2 and 1-by-1 mip levels), then
returns the expected RGBA image. Extra compressed bytes are accepted and discarded after the
declared level has been consumed; a too-short DXT level throws \texttt{std::out\_of\_range}, not
a unified content-load error.

The uncompressed branch has a sharper current-develop defect. It calls
\texttt{ReadBytesExactOrThrow()}, which proves only that the file supplied as many bytes as its
own \emph{declared} count, then indexes the result as though it were exactly
\texttt{faceSize * faceSize * 4}. A small positive raw byte count therefore reaches
\texttt{bytes[o]} through \texttt{bytes[o+3]} past the vector's end. This is not merely a missing
test assertion: the later \texttt{feature/audit} commit \texttt{e3bc2be1} records it as a
reproduced heap out-of-bounds read and adds the sibling readers' equality guard and regression
test, but that commit is not an ancestor of current \texttt{develop}. A too-large raw count is
silently ignored after the expected colors have been formed. Treat both cases as the live source
contract until the source fix is integrated; the earlier broad XNB hardening commit's description
does not override the code actually present on this branch.

The preflight before those levels is only partial. It rejects a nonpositive size and presents
\texttt{size * size * 4} to \texttt{CheckDecodedByteSize()}, bounding one base face's nominal
RGBA bytes to the default 256 MiB. Six faces and a complete mip chain can require almost eight
times that base amount, and there is no total cube-data budget. The raw signed
\texttt{int64\_t size * size * 4} expression can itself overflow before the helper observes a
value; its checked-multiply counterpart is also only on unmerged \texttt{feature/audit}. There
is no reader-level device-size admission either (apart from a separate D3D9 profile check inside
the resource constructor). The base face cap is consequently neither a complete arithmetic guard
nor a portable GPU limit.

Mip handling repeats the family pattern. The constructor sees only \texttt{levels > 1}, then
creates the framework's full square chain rather than a chain exactly equal to the file count.
Zero or negative levels create a one-level cube and read no face payload; too many positive levels
continue with 1-by-1 data and are never rejected. TextureCube validates a nonnegative level, face,
rectangle, and element count, but not that the level is allocated. Its mip helper shifts the base
size by the supplied level, so sufficiently large hostile levels make that shift undefined before
the renderer update. The reader also has no generic cube-capability gate: a renderer using the
default null cube factory produces an object whose upload/readback methods simply skip renderer
work. A successful parser return is not proof that an actual cube resource exists.

Existing-instance behavior is particularly unlike FNA's reference object. Normal ContentReader
dispatch gives this reader no instance. For a direct move-only \cnaclass{TextureCube} read,
CNA's type erasure requires the caller to supply a \texttt{shared\_ptr<TextureCube>}; it moves the
pointed-to resource into a local value, uploads into that value without checking its size, format,
device, or mip chain against the file, then returns a newly boxed result. The caller's original
shared pointer is left pointing at a moved-from resource. FNA instead writes directly into its
non-null reference instance. A direct existing-instance read therefore has neither FNA identity
semantics nor a safe replacement/matching contract in CNA.

This reader has the strongest external fixture among the three texture variants, but still a
limited oracle. A real MonoGame \texttt{SampleCube64DXT1Mips.xnb} loads through
\cnaclass{ContentManager}; it covers all six faces and seven DXT1 levels from 64 through 1,
including the sub-4-by-4 block case. The test checks size, Color output format, level count,
non-uniform PositiveX level 0, and that one NegativeZ 1-by-1 read does not throw. It does not
check exact pixels, face orientation/order, every face, DXT3/DXT5, raw Color, bad byte count,
size/level/format/version/device/existing cases, cache/unload, or another producer/platform.
There is no Cube seed in the whole-container fuzzer. Thus the fixture is valuable evidence for
ordinary DXT1 traversal, not proof that arbitrary Cube XNB input is bounded or that all six faces
are semantically correct.

\subsection{SpriteFontReader has real DXT3 and LZX fixtures, but does not validate its parallel glyph tables}
\label{sec:xnb-spritefont-reader-contract}

The explicitly registered
\texttt{Microsoft.Xna.Framework.Content.SpriteFontReader} has FNA's ordinary fresh-font payload,
but it is an object graph rather than a run of bare fields. In this exact order it dispatches a
\cnaclass{Texture2D} atlas, a \texttt{List<Rectangle>} of glyph bounds, a second
\texttt{List<Rectangle>} of cropping rectangles, a \texttt{List<Char>} character map, Int32
\texttt{lineSpacing}, Single \texttt{spacing}, a \texttt{List<Vector3>} of kerning triples, and
a Boolean followed by a Char only when that Boolean says a default character is present. Each
nested object begins with its own one-based XNB type-reader index. The final optional character
is therefore a literal Boolean/Char pair, not a \texttt{NullableReader<Char>} object. CNA applies
its exact serialized-version-zero rule to the table entries around this FNA-shaped body.

The reader's registration has a deliberately narrow dependency surface. It installs its three
actual closed generic names --- ListReader for Rectangle, Char, and Vector3 --- but intentionally
does not install the atlas's Texture2DReader; a standalone registration caller must do that
separately. The built-in umbrella does install both. CNA's generic-list adapter recreates its
element reader from the global canonical-name registry, rather than retaining FNA's matching
instance from the current file's reader table. Thus only these three closed combinations are
available; a dependency spelling CNA has not registered fails rather than being reflectively
supplied. The real standard table has exactly the eight names needed by this route.

Each list first reads a signed Int32 count and the common collection guard accepts only zero
through the configured ten-million-element limit. That protects an individual reservation, but
there is no aggregate budget across the three lists and no semantic relation between them. The
constructor simply stores all four parallel vectors and builds its character lookup by assigning
each map entry its last index. It does not require equal glyph/cropping/character/kerning counts,
sorted or unique characters, finite kerning and spacing, positive line spacing, or rectangles
inside the atlas. A malformed font can consequently load successfully and later make
\texttt{MeasureString()} or \texttt{SpriteBatch::DrawString()} index a shorter kerning, cropping,
or glyph vector out of bounds. FNA also trusts its content-pipeline tables, but its managed lists
fail with a bounds exception; CNA's unchecked \texttt{std::vector} indexing is not a safe error
boundary for untrusted XNB input.

The optional default character exposes a second current-develop defect. CNA copies a present
character without checking that the character map contains it. When an unknown text character
then falls back, both measurement and drawing call \texttt{find(defaultCharacter)} and immediately
read \texttt{it->second}; an absent default therefore dereferences \texttt{end()}. FNA's
dictionary indexer instead throws on that inconsistent state. The later unmerged
\texttt{feature/audit} commit \texttt{fb1dcbc4} adds construction/setter validation and a checked
fallback, but it is not an ancestor of current \texttt{develop}. This is observable in the
current CNJ SpriteFont test fixture too: it supplies \texttt{"?"} as the default while its sole
glyph is \texttt{'A'}, then never measures an unknown character. Treat the absence of an XNB
default-character invariant as live source behavior, not as a harmless test omission.

Existing-instance handling differs deliberately, even though ordinary manager dispatch never
reaches it. SpriteFontReader leaves \texttt{CanDeserializeIntoExistingObject} false, as FNA
does, so normal \cnaclass{ContentManager} loading creates a fresh value. A direct advanced read
with an existing C\texttt{++} font instead throws \texttt{System::NotImplementedException} before
consuming payload. FNA's otherwise-dead branch reloads only the existing atlas texture in place,
discards every remaining font field, and returns that same instance. CNA consequently has no
direct reload equivalence, but it also avoids inventing a public mutator solely for a normal-path
case that cannot occur.

The evidence is appreciably stronger than for Texture3D but bounded. Real uncompressed MonoGame
\texttt{Default.xnb} goes through \cnaclass{ContentManager} with the standard eight-reader table:
its 128-by-128 DXT3 atlas, all four 95-element lists, characters from space through tilde,
line spacing 19, zero extra spacing, and no default character are independently recorded in its
fixture manifest. The test asserts the scalar and character-range subset, not glyph, cropping, or
kerning values, decoded atlas pixels, or rendering. A second real MonoGame
\texttt{FontCalibri14.xnb} exercises the same reader after a 44,032-byte multi-block LZX decode;
the test establishes only a nonempty map and positive \texttt{MeasureString("Hello")} width.
The focused reader tests cover registration and the false existing-instance property, while the
many renderer SpriteFont pixel tests hand-build tiny fonts rather than validating XNB tables. No
located test covers a true XNB default character, list/count mismatch, duplicate or malformed
character map, invalid rectangles/floats, truncation, nonzero reader version, direct-existing
failure, cache/unload, another producer/platform, or a SpriteFont seed in the container fuzzer.
Treat the two real fixtures as sound ordinary-layout and compressed-path evidence, not as a
validation or adversarial-input oracle.

\subsection{ModelReader reconstructs the ordinary shared-resource graph, but accepts contradictory hierarchy and draw metadata}
\label{sec:xnb-model-reader-contract}

The registered \texttt{Microsoft.Xna.Framework.Content.ModelReader} depends on three other
registered readers: \texttt{VertexDeclarationReader}, \texttt{VertexBufferReader}, and
\texttt{IndexBufferReader}.  A declaration is an Int32 stride, an Int32 element count, then an
Int32 offset, format, usage, and usage index for every element.  A vertex buffer embeds that
declaration through a raw-reader call, then has a UInt32 vertex count and exactly
\texttt{vertexCount * stride} raw bytes.  An index buffer is a Boolean selecting 16- or 32-bit
indices, an Int32 byte count, and that many bytes.  CNA limits declaration elements to 1,024,
rejects a vertex count which cannot fit its signed resource API, widens the vertex-byte
multiplication before the exact read, and rejects a negative stride or overflowing byte total.
Those are useful hardening additions around FNA's direct reads.  It does not validate declaration
offsets, enum values, or their relationship to the stride, however.  Index data needs only be
nonnegative: a non-multiple byte count is accepted, divided down to an index count, and its one to
three trailing bytes are silently discarded.  The reader passes the byte vector on by a typed
pointer cast rather than first requiring an exact element-width multiple.

At the start of the root reader CNA requires both an owning content manager and its graphics
device; the result owns live vertex/index buffers and effects, not just a portable serialized
description.  It then reads a UInt32 bone count (capped at 100,000), each bone's XNB String and
Matrix, and each hierarchy record.  A bone reference is a one-based ID with zero for null: it
occupies one byte while the model has fewer than 255 bones, and UInt32
otherwise.  Each record contains a parent reference followed by a UInt32 child count and that
many child references.  Mesh count and each mesh-part count are signed and independently limited
to the same 100,000 ceiling; mesh records carry an XNB String name, a bone reference, a
\cnaclass{BoundingSphere}, Tag object, and parts.  A part has four signed draw scalars in the
order vertex offset, vertex count, start index, and primitive count, its Tag, then shared
VertexBuffer, IndexBuffer, and Effect reference indices.  The final root-bone reference and
model Tag finish the body.  This is FNA's version-zero wire order, with CNA's normal reader-table
version admission around it.

The hierarchy fields are redundant in a well-formed pipeline output, but CNA does not use them
symmetrically.  It reads and completely ignores each scalar parent reference; only processing a
parent's child list calls \texttt{ModelBone::AddChild()}, which sets the child's parent pointer.
FNA assigns a non-null scalar parent \emph{and} adds children.  CNA therefore accepts an
out-of-range scalar parent and accepts mutually inconsistent parent/child declarations whenever
the child references themselves are in range.  There is no child-count cap, duplicate-child,
single-parent, cycle, or parent-before-child check.  A malformed graph can consequently make
\texttt{CopyAbsoluteBoneTransformsTo()} compose a bone with a destination transform that has not
yet been calculated; \cnaclass{Model}'s linear index-order loop was written for ordinary
pipeline hierarchy ordering, not graph repair.  CNA does validate a non-null mesh parent and a
non-null root reference.  But a null root becomes index zero when bones exist (and null only when
the bone vector is empty), whereas FNA immediately indexes its \texttt{-1} result and fails.

Mesh and draw metadata are similarly trusted after their coarse counts.  CNA preserves the mesh
name and bounding sphere and assigns the validated mesh parent through Model's ownership-aware
constructor, but does not check that offsets, vertex counts, starts, or primitive counts are
nonnegative or fit their resolved buffers.  A positive shared-resource index is range-checked by
the generic reader and fixed up only after every shared object has been read; index zero is a
legal null reference, leaving that part's corresponding pointer unset.  No ModelReader check
requires the three pointers before later rendering.  A positive primitive count with an effect
and absent buffer is thus not rejected at loading time.  The normal two-pass ordering is
important: callbacks install raw buffer/effect pointers into the already built parts only after
the root Model exists, while a private \texttt{shared\_ptr}-owned resource bundle retains the
bones, meshes, parts, buffers, and effects for every copy of that Model.  This is why cache copies
can safely share the same GPU resources rather than re-decoding them.

Tags and existing-instance reads are explicit compatibility boundaries rather than unnoticed
omissions.  CNA consumes every model, mesh, and part Tag so the stream remains aligned, but throws
\cnaclass{ContentLoadException} for any non-null value.  FNA stores those objects on the model
objects.  The default false \texttt{CanDeserializeIntoExistingObject} flag means a normal
manager load has no existing Model in either implementation.  A direct CNA read with one instead
throws before consuming the body; FNA's otherwise unreachable branch reuses its existing mesh
parts and returns the same Model after consuming the remaining root and Tag.  Normal CNA loading
uses the copyable generic asset cache, so successive \texttt{Load<Model>()} results share the
reader-owned graph; it is not an in-place reload API.

The evidence is good for the intended ordinary case but narrow for this graph's edge cases.  A
real externally produced, uncompressed MonoGame version-five \texttt{BlenderDefaultCube.xnb}
passes through the complete manager path: two bones (\texttt{RootNode} then \texttt{Cube}), one
Cube mesh parented to the second bone, its radius-$\sqrt{3}$ bounding sphere, one 24-vertex/
12-primitive part, a 24-byte Position--Normal declaration, a 16-bit 36-index buffer, and a
\texttt{BasicEffect}.  The focused test also proves that two manager loads retain pointer identity
for the vertex buffer.  The deterministic whole-container fuzzer mutates that complete fixture
1,500 times, exercising root dispatch and shared-resource fixups; it is a valuable clean-failure
screen, not a semantic graph test.  There is no compressed or non-MonoGame Model fixture, no
fixture with two parts sharing a resource, no 255-bone/UInt32-reference case, and no targeted
test for Tags, existing models, null or invalid shared references, empty/null roots, contradictory
or cyclic hierarchy, out-of-range draw scalars, non-multiple index data, ownership after unload,
or a model draw on a real renderer.  Successful loading of the cube therefore establishes the
standard pipeline path, not a promise that arbitrary Model XNB metadata is safe to draw.

\subsection{Stock-effect readers reproduce their material fields, but external textures and four of five formats lack real-asset evidence}
\label{sec:xnb-stock-effect-reader-contract}

The five registered stock-effect readers are a deliberately separate family from the general
\texttt{Microsoft.Xna.Framework.Content.EffectReader}.  The latter is the recognised-but-
unsupported compiled platform-shader route described earlier; it consumes no effect bytecode and
throws a named content-load error.  The five concrete readers instead construct CNA's corresponding
native stock effects.  They all erase their C\texttt{++} result to \texttt{shared\_ptr<Effect>}.
This is not a claim that all five have the same runtime
class: it is the necessary common type for a Model mesh part's polymorphic shared-effect callback.
An \texttt{any} cast in this implementation requires an exact C\texttt{++} type, whereas FNA's
managed reference cast can naturally upcast any of the five concrete objects to \texttt{Effect}.

Each body starts by constructing its effect from the owning content manager's GraphicsDevice, so
even an entirely texture-free stock-effect XNB cannot be read standalone or by a manager without a
device.  Its fields then have FNA's literal version-zero order.  BasicEffect has an external
Texture2D reference; diffuse, emissive, and specular Vector3 values; Single specular power and
alpha; and Boolean vertex color.  AlphaTestEffect has an external Texture2D; Int32 comparison
function; UInt32 reference alpha; diffuse Vector3; Single alpha; and Boolean vertex color.
DualTextureEffect has two external Texture2D references, diffuse Vector3, Single alpha, and
Boolean vertex color.  EnvironmentMapEffect has external Texture2D and TextureCube references,
then Single map amount, Vector3 map specular, Single Fresnel factor, diffuse and emissive Vector3
values, and Single alpha.  SkinnedEffect has external Texture2D, Int32 weights per vertex,
diffuse/emissive/specular Vector3 values, Single specular power, and Single alpha.  Fog, world/
view/projection matrices, lights, and the other runtime effect switches are not serialized in this
family; each remains at its native constructor default.

The texture strings use the generic \texttt{ReadExternalReference<T>()} path, not the Song/Video
filesystem shortcut.  Thus an empty string is null, while a nonempty relative reference resolves
beside the effect asset, goes through normal ContentManager resolution, and retains that returned
texture through a private \texttt{shared\_ptr}.  CNA applies the generic lexical escape rejection
for a path which normalizes above the root, although the existing generic absolute-path and
symlink limits still apply.  It also asks directly for \texttt{Texture2D} or \texttt{TextureCube}.
FNA asks for the broader \texttt{Texture} type and uses \texttt{as Texture2D} or
\texttt{as TextureCube}; a wrong concrete root therefore becomes null in FNA but reaches CNA's
exact-type loading/cast boundary.  No stock-effect test exercises a nonempty, missing, escaping,
wrong-type, absolute, or symlinked external texture reference.

Apart from the normal scalar reads, the readers intentionally leave validation to the effect
properties.  SkinnedEffect's setter admits only one, two, or four weights per vertex, matching
FNA's own property rule, so another Int32 fails while the reader is constructing the effect.
The other material floats and Vector3 components are not required to be finite or inside a color/
opacity range.  AlphaTestEffect also casts any Int32 to \cnaclass{CompareFunction} without an enum
check and casts the full UInt32 reference-alpha field into CNA's signed Int32 property.  This
mirrors FNA's raw casts rather than enforcing the public property's nominal 0--255 range.  A
well-formed pipeline normally supplies sensible values; a successful reader return is not a
sanitisation pass for arbitrary material metadata.

Existing-instance behavior has a less surprising C\texttt{++} ownership boundary than the
move-only audio readers.  These reader targets are copyable \texttt{shared\_ptr<Effect>} values,
but their default false \texttt{CanDeserializeIntoExistingObject} capability means normal root
dispatch supplies none.  A direct caller must box the common base pointer rather than one typed
to a concrete effect.  CNA unboxes then ignores it and returns a fresh effect, as FNA also ignores
its dormant existing parameter.  The caller's original shared pointer remains
valid.  Within a Model, the delayed shared-resource callback installs a raw effect pointer in the
part and the Model's owner bundle retains the shared pointer, so repeated Model loads preserve the
same live material resources.

Evidence divides sharply between BasicEffect and the other four readers.  The five canonical
factory names and their hand-written field orders have focused registration/direct-reader tests.
The real externally produced MonoGame \texttt{BlenderDefaultCube.xnb} provides a 46-byte
BasicEffect shared-resource body with no texture, the recorded diffuse/specular material values,
and a full Model integration route; the Model fuzzer mutates that same BasicEffect path as part of
its 1,500 whole-container cases.  AlphaTest, DualTexture, EnvironmentMap, and Skinned tests are
hand-constructed empty-texture streams checked field by field against current FNA sources.  No
reader has a standalone root XNB manager load, nonzero reader-version or truncation test,
external-texture success/failure test, direct-existing test, cache/unload check, or its own
fuzzer seed; no real asset was located for the latter four.  Treat current support as a credible
ordinary material-field port, not evidence that every stock-effect XNB producer, texture graph,
or malformed payload is interoperable.

\subsection{SoundEffectReader widens its real WAV support, but its wrapped-loop endpoint can overflow}
\label{sec:xnb-soundeffect-reader-contract}

The registered \texttt{Microsoft.Xna.Framework.Content.SoundEffectReader} first reads a UInt32
\texttt{formatLength}, then the six WAVEFORMATEX fields: UInt16 format tag and channel count,
UInt32 sample rate and average bytes per second, and UInt16 block alignment and bits per sample.
For a declared length above 16 it consumes a swapped UInt16 \texttt{cbSize}; an XMA2/34-byte
extension is field-read solely to preserve stream position, while another extension is retained
as raw bytes. Next come a signed Int32 data length and exact data bytes, signed Int32 loop start
and length, and a UInt32 stored duration in milliseconds. On Xbox platform \texttt{'x'} only the
WAVEFORMATEX multi-byte fields are byte-swapped; the outer XNB counts, loop values, and audio data
remain little-endian. This matches FNA's unusual swap boundary. CNA additionally requires the
usual serialized reader version zero; FNA's reader table does not enforce that version.

The length word is an account of the format block, not a full structural admission check. CNA
unconditionally consumes the 16-byte base fields before consulting it, so an underdeclared value
at or below 16 is not independently rejected. Above 16, a length too small for \texttt{cbSize} or
too large to represent as the reader's signed exact-read count fails cleanly; the focused tests
cover the ordinary 16/18/extended sizes and adversarial extremes. The raw audio count likewise
uses \texttt{ReadBytesExactOrThrow()}, so negative declarations become a content-load error and a
truncated stream becomes end-of-stream rather than the old pointer/count overread. It has no
separate decoded-audio or aggregate asset-size budget, however: a physically supplied, valid but
very large audio body may still be retained and handed to a decoder.

Current format support is wider than several stale local artifacts claim. Six real MonoGame
fixtures and their direct-reader tests establish 16-bit PCM mono/stereo, 8-bit PCM, 32-bit IEEE
float, MS-ADPCM, and IMA-ADPCM. PCM16 takes CNA's raw signed-little-endian fast path. The other
four formats are put in a synthetic WAV and decoded by SDL3 through \texttt{SoundEffect::FromStream};
for real MonoGame MS-ADPCM whose \texttt{cbSize} is zero, CNA synthesizes the standard seven
coefficient pairs and a samples-per-block value before calling SDL. Channels other than mono or
stereo, unknown tag/bit combinations, and XMA2 are rejected with an asset-specific
\texttt{ContentLoadException}; CNA parses the XMA2 extension only to remain aligned, while FNA
forwards that form to its own SoundEffect/FAudio path. Thus XMA2 is a genuine CNA format gap, not
an unrecognized reader name.

The source, header, and tests are authoritative here: the older XNB support page and the four
PCM8/float/MS-ADPCM/IMA-ADPCM fixture manifests still label those working formats ``rejected.''
Their status predates the WAV-wrapper change and must not be copied into a current compatibility
claim. The stored duration also has a CNA-only role. FNA reads then discards it; CNA compares a
nonzero declaration with the decoded duration and rejects a ratio below 0.5 or above 2.0. That
generous check catches gross sample-rate/channel errors and accepts observed ADPCM block rounding,
but it can reject a malformed or unusual asset that FNA would have accepted.

Loop values are nominally decoded sample-frame values and are deliberately not bounded against
the decoded audio length, matching FNA's internal constructor. The direct PCM16 route stores the
two signed fields after unsigned conversion, so negative values wrap as they do in FNA. The
WAV-wrapper route adds a CNA-specific hazard: for a positive \texttt{loopLength},
\texttt{AppendSmplChunkIfLooped()} computes signed \texttt{loopStart + loopLength} before writing
the WAV \texttt{smpl} end field. An extreme positive pair can overflow that signed Int32 addition
(undefined behavior) before the intended unsigned encoding. A normal IMA-ADPCM test proves
1500/400 frame values survive a roughly 4:1 compressed payload and a loopless test proves zeroes
stay absent, but no test reaches the overflow or validates loop bounds. Treat malformed wrapped
loops as an outstanding current-develop safety gap rather than relying on the direct PCM policy.

The reader does not deserialize into an existing object in normal manager dispatch, as FNA also
advertises. There is a C\texttt{++}-specific direct-path consequence because \cnaclass{SoundEffect}
is move-only: type erasure accepts a caller's \texttt{shared\_ptr<SoundEffect>} and moves its
pointed-to value into the optional existing argument, but SoundEffectReader ignores that argument
and returns a new effect. The caller's pointer consequently refers to a moved-from object. FNA
also ignores its supplied instance, but its reference object remains intact while it returns a
new one. Normal \cnaclass{ContentManager} loads avoid this path and create a fresh move-only effect
on every call rather than using the generic copyable-object cache.

The test matrix is substantial but not a universal compatibility proof. One manager-level test
loads real 16-bit MonoGame PCM and asserts its name and positive duration. Six focused tests load
all real supported variants through \texttt{ReadAsset}, and a hand-built XMA2 body proves the
rejection branch. Further focused cases cover Xbox byte swapping, malformed data/format lengths,
duration-oracle decisions, mono/stereo admission, normal loops, and decoder error context; a
deterministic 1,120-combination WAVEFORMATEX boundary sweep accepts only a result or clean
content/end-of-stream error. The whole-container fuzzer mutates the real PCM16 asset 1,500 times.
Missing are a real XMA2 or Xbox asset, an external FNA sample-level decode comparison, a true
LZX-compressed audio XNB (the misleadingly named \texttt{Explosion.xnb} is a Texture2D), extreme
loop tests, large valid-body limits, direct-existing behavior, and non-MonoGame producers. The
reader is therefore well exercised for current supported desktop WAV paths, but neither every
upstream codec nor every hostile metadata combination is established.

\section{Registration and platform-dependent inventory}

\texttt{RegisterAllBuiltInXnbReaders()} is the single aggregate registration point and is
idempotent. It is deliberately not invoked by the \cnaclass{ContentManager} constructor,
because tests and embedders can clear and repopulate the process-wide factory registry. A real
application that expects stock XNB assets must ensure registration occurs before the first load.

The resulting inventory is platform-dependent. Video is compiled and registered only when
FFmpeg is available; decimal support also depends on the required integer representation. A
fixed reader count without those build conditions is therefore misleading.

General compiled \texttt{EffectReader} input is registered as a deliberate poison pill. Its
factory creates a reader whose failure names compiled platform shader bytecode as the unsupported
reason, producing a specific diagnostic instead of the vague ``unknown reader'' that omission
would cause. The five stock-effect readers remain separate supported formats.

\section{Object-graph rules worth carrying into custom readers}

Custom binary readers participate in the same session rules as built-ins:

\begin{enumerate}
  \item type-reader indices are one-based, with zero representing null;
  \item reader versions currently require exact version zero because no reader widens
        \texttt{SupportsVersion};
  \item nested object depth is bounded and restored by an RAII guard on every exit path;
  \item shared objects are all read before queued fixups run;
  \item external references are confined as file-internal references even though the caller's
        top-level absolute asset name may be allowed by \cnaclass{ContentManager};
  \item move-only values are boxed through \texttt{shared\_ptr<T>} because \texttt{std::any}
        itself requires a copy-constructible payload.
\end{enumerate}

These are not implementation trivia. A custom reader that bypasses the session or invents a
parallel registry can lose null semantics, resource identity, bounds, or error context even when
its field decoding is individually correct.
